\section{Methodology}
\label{sec:methodology}

\subsection{Data source and filtering}
The underlying data is OpenAlex's public authors snapshot~\citep{openalex}, distributed as Parquet on \texttt{s3://openalex/data/parquet/authors/} (\texttt{us-east-1}, no credentials required) and partitioned into over 1,600 daily \texttt{updated\_date=} directories totalling about 53\,GB. Only four columns are read from each remote file (\texttt{id}, \texttt{h\_index}, \texttt{affiliations}, \texttt{topics}), and an author is retained only if all of the following hold:
\begin{itemize}
	\item \texttt{summary\_stats.h\_index} is non-null and greater than zero;
	\item at least one entry in \texttt{affiliations} has \texttt{institution.type = 'education'};
	\item at least one entry in \texttt{topics} carries a non-null field classification.
\end{itemize}
Applying these filters to the full snapshot yields 17.1 million authors, spanning 20,669 institutions and OpenAlex's current list of 26 fields. One limitation carried over from the source data: an author's \texttt{h\_index} is computed by OpenAlex from all works meeting its citation criteria regardless of retraction status, which OpenAlex records inconsistently~\citep{hauschke2025}; a small, unquantified fraction of h\textsubscript{1} values in the dataset may therefore be inflated by since-retracted work.

\subsection{Author-level assignment}
\label{subsec:author-assignment}
Each retained author is reduced to a single row: one institution and one field. An author may list several affiliations with \texttt{institution.type = 'education'}; rather than fractionally splitting or duplicating that author across institutions, I assign them to the one with which they were most recently affiliated. Recency is read from OpenAlex's \texttt{affiliations} field, which records, for each institution an author has ever listed, the calendar years in which they published while there; I take the institution whose most recent such year is latest, breaking ties (identical latest year) by the numerically smallest OpenAlex institution identifier. I deliberately do not use the similarly-named \texttt{last\_known\_institutions} field for this: despite its name, it is simply the set of affiliations listed on an author's single most recent work, carrying no per-entry recency of its own, so when an author has more than one such affiliation (e.g., a joint appointment or co-affiliated coauthors misattributed to one author) there is no way to tell which is ``more recent'' than another.

Field assignment works the same way but is weighted. OpenAlex tags each author with a set of topics, each carrying a \texttt{count} (roughly, the number of the author's works classified under that topic) and a parent field. I sum \texttt{count} within each field and assign the author to whichever field carries the largest total weight: their modal field by publication volume, not simply their single most granular topic. The result is an \texttt{authors} table with one row per author: \texttt{(author\_id, h\_index, institution\_id, institution\_name, field, field\_name)}.

\textcolor{red}{Two structural features of this procedure limit the impact of field classification errors. First, the paper operates at the \emph{field} level (26 fields) rather than at the finer topic level (${\sim}4{,}500$ topics); OpenAlex's topic classifier fine-tunes a multilingual BERT model on title, abstract, and citation data and achieves a top-1 accuracy of 0.53~\citep{vaneck2024}, but many topic-level errors stay within the same parent field, so field-level accuracy is substantially higher. Second, an author's field is determined by the mode of their \emph{entire publication record}, not by a single work: a researcher would need a substantial fraction of their career output systematically misclassified to a different field before their assignment would change.}

\subsection{The h-index and its successive generalization}
Hirsch's~\citeyearpar{hindex} original index, here $\mathrm{h}_1$, is computed by OpenAlex per author from citation counts and taken as given. Schubert's generalization~\citeyearpar{parent} lifts the same definition up a level of aggregation: a group of $N$ members, each already carrying an $\mathrm{h}_{k-1}$ index, has an index $\mathrm{h}_k = n$ if $n$ of its members have $\mathrm{h}_{k-1} \geq n$ and the remaining $N-n$ members have $\mathrm{h}_{k-1} < n$:

\begin{algorithm}
\caption{Successive h-index of a group $G$, given a value $v(m)$ for each member $m \in G$}
\begin{algorithmic}[1]
\State Sort $G$ by $v$ descending; let $r(m)$ be $m$'s 1-indexed rank in that order
\State \Return $\max \{\, r(m) : v(m) \geq r(m),\ m \in G \,\}$
\end{algorithmic}
\end{algorithm}

I use this rank-based procedure, implemented as a window function in SQL ranking rows within each group and taking the largest rank whose value has not yet fallen below it, for every successive index in this paper: $\mathrm{h}_2$ of an institution over the $\mathrm{h}_1$ of its authors, and $\mathrm{h}_3$ of a country over the $\mathrm{h}_2$ of its institutions.

\subsection{h\textsubscript{2}: institution and institution$\times$field}
h\textsubscript{2} is computed twice, over different groupings of the same author-level data. The field-specific table groups authors by (institution, field) and yields 288,300 pairs; institutions are matched on their numeric identifier only; institution and field \emph{names} are resolved afterward via \texttt{arg\_max}, so that spelling or formatting variants of the same name across snapshot partitions cannot silently split one institution into several groups. The institution-overall table ignores field entirely and applies the same procedure across all of an institution's authors regardless of discipline, yielding one h\textsubscript{2} per institution (20,669 rows).

\subsection{h\textsubscript{3}: country level}
Institutions are mapped to countries using the \texttt{country\_code} field of OpenAlex's institutions snapshot (also read directly from S3, joined on institution ID). h\textsubscript{3} is then computed with the same rank-based procedure, once globally per country over the h\textsubscript{2} values of its institutions and once per (country, field) over the field-specific h\textsubscript{2} values, producing one file per field.

\subsection{Efficiency}
\label{subsec:methods-efficiency}
\cite{mryglod2022} showed that group-level h-indices grow with group size even when per-member quality is held constant. They were not studying successive h indices, but a similar argument makes raw $\mathrm{h}_2$ and $\mathrm{h}_3$ values unsuitable for comparing institutions or countries of unequal scale.

\cite{egghe2008} applied a Lotkaian framework to successive h indices. Suppose the article--citation frequency function follows Lotka's law with exponent $\alpha_0 > 1$, and the author--article frequency function follows Lotka's law with exponent $\alpha_1 > 1$. Egghe showed that the distribution of $\mathrm{h}_1$ values across authors is itself Lotkaian with compound exponent $\beta_1 = \alpha_0\alpha_1$. Applying the h-index formula for Lotkaian systems, the institution-level index satisfies
\begin{equation}
	\mathrm{h}_2 = \mathrm{S}^{1/\beta_1}
	\label{eq:h2-scaling}
\end{equation}
where S is the number of authors at the institution. An analogous argument at the next level gives the distribution of $\mathrm{h}_2$ values across institutions a Lotkaian exponent $\beta_2 = \alpha_0\alpha_1\alpha_2$, where $\alpha_2$ is the Lotka exponent of the institution--author frequency function, so that
\begin{equation}
	\mathrm{h}_3 = \mathrm{R}^{1/\beta_2}
	\label{eq:h3-scaling}
\end{equation}
where R is the number of institutions in the country. I define size-normalized efficiency metrics by dividing out these predicted scalings:
\begin{equation}
	\varepsilon_2 = \frac{\mathrm{h}_2}{\mathrm{S}^{1/\beta_1}}
	\label{eq:eff2}
\end{equation}
\begin{equation}
	\varepsilon_3 = \frac{\mathrm{h}_3}{\mathrm{R}^{1/\beta_2}}
	\label{eq:eff3}
\end{equation}
An institution with $\varepsilon_2 > 1$ outperforms the Lotkaian prediction for its headcount; one with $\varepsilon_2 < 1$ underperforms it.

The exponents $\beta_1$ and $\beta_2$ are estimated directly from the data rather than assembled from separately fit $\alpha_0$, $\alpha_1$, $\alpha_2$. I apply discrete maximum-likelihood tail estimation following \citet{clauset2009}, with $x_\text{min}$ chosen to minimize the Kolmogorov--Smirnov distance between the empirical and fitted CDFs, to the distributions of $\mathrm{h}_1$ across all authors and $\mathrm{h}_2$ across all institutions, reading $\beta_1$ and $\beta_2$ as their fitted Lotka exponents. I then fit $\alpha_0$, $\alpha_1$, and $\alpha_2$ individually: $\alpha_0$ from a stratified sample of per-article citation counts drawn from the OpenAlex works snapshot, $\alpha_1$ from per-author work counts, and $\alpha_2$ from per-institution author counts. Finally, I verify that $\alpha_0\alpha_1 \approx \beta_1$ and $\alpha_0\alpha_1\alpha_2 \approx \beta_2$ as a cross-check that the compounding-exponent structure of Egghe's model holds in the OpenAlex dataset.

\subsection{Supplementary metrics}
\label{subsec:methods-supplementary}
Several secondary statistics are derived from the h\textsubscript{2} and h\textsubscript{3} tables to probe what the headline numbers obscure:
\begin{itemize}
	\item \textbf{Gini coefficient}, computed on the within-field distribution of institutional h\textsubscript{2}, quantifies whether a field's excellence is concentrated in one institution (Gini near 1) or broadly distributed (Gini near 0).
	\item \textbf{Breadth score} counts, per institution, how many of the 26 fields it ranks in the global top 10, top 50, and top 100.
	\item \textbf{Non-biomedical ranking} recomputes institution-level h\textsubscript{2} after excluding authors whose modal field is Medicine, Biochemistry/Genetics/Molecular Biology, Immunology and Microbiology, Neuroscience, or Health Professions, to test whether the overall ranking is, in substance, a ranking of medical school size.
\end{itemize}

\subsection{Manual cross-validation}
\label{subsec:methods-crossval}
Because errors in OpenAlex's underlying author disambiguation or citation counts would propagate directly into every successive index, I manually spot-check a targeted, non-random sample of authors rather than a random one, chosen at the points where an error would be both most likely and most consequential: authors sitting exactly at an institution's h\textsubscript{2} threshold, where a $\pm 1$ error in a single h\textsubscript{1} could change the institution's rank; top authors in fields where OpenAlex's topic classification is known to be thinner (e.g., Arts and Humanities); top authors at institutions with non-Anglophone naming conventions, where author-name disambiguation is more error-prone and where OpenAlex's coverage is known to diverge unevenly from other databases by country, most severely for China~\citep{zheng2025}, including institutions suspected of affiliation-swapping for grant purposes~\citep{saudi}; and the small number of authors with the highest h\textsubscript{1} globally, as a sanity check at the extreme of the distribution. For each candidate, I resolve the author's OpenAlex profile through the public API and compare it by hand against Web of Science, Scopus, and Google Scholar. \textcolor{red}{This design is deliberately adversarial: it is built to surface errors, not to estimate how frequently they occur across the full population, so any outlier rate found in it bounds the qualitative failure mode (which authors, and by how much) rather than the population-wide error rate.}

\subsection{\textcolor{red}{Sensitivity to disambiguation error}}
\label{subsec:methods-sensitivity}
\textcolor{red}{Because the cross-validation sample above cannot itself bound a population-wide error rate, I separately ask how much author-disambiguation error the paper's conclusions could tolerate before changing using a Monte Carlo sensitivity analysis run directly on the full authors table. At a population-wide error rate $p$, each author is independently flagged with probability $p$; a flagged author's $\mathrm{h}_1$ is multiplied by an inflation ratio drawn uniformly from the four ratios actually observed among the cross-validation outliers, simulating a merged-identity error of the same kind and magnitude found by hand. I recompute institution-level h\textsubscript{2}, country-level h\textsubscript{3}, and the field-specific h\textsubscript{2} rankings, and compare each to the unperturbed baseline via Spearman rank correlation and top-set overlap. Rather than picking a single rate to call a ``worst case,'' I run five independent trials at each of six rates, $p \in \{1\%, 5\%, 10\%, 20\%, 40\%, 80\%\}$, tracing out how ranking stability responds to the assumed amount of error from a plausible rate up to a deliberately implausible one, so that the paper's tolerance for this kind of error is characterized by a curve rather than asserted at a single point.}

\subsection{\textcolor{red}{Robustness to field assignment}}
\label{subsec:methods-field-robustness}
\textcolor{red}{The single-field assignment described in Subsection~\ref{subsec:author-assignment} asks ``which field accounts for the largest share of an author's publication volume?'' and assigns each author exclusively to that field. To test whether field-specific h\textsubscript{2} rankings are stable under an alternative that allows multidisciplinary authors to count in multiple fields simultaneously, I re-run the field-specific h\textsubscript{2} computation after expanding each author's field membership: at a threshold $\tau$, an author is assigned to every field in which their OpenAlex topic-weight share is at least $\tau$ (where topic-weight is the same count-weighted field sum used in the original assignment). I test two thresholds, $\tau \in \{20\%, 33\%\}$, spanning from liberal (an author need contribute only one fifth of their output to count in a field) to moderate (one third or more). For each threshold I recompute h\textsubscript{2} for every (institution, field) pair in the expanded author roster, using the same rank-based algorithm as the original, and compare to the baseline via Spearman rank correlation on shared (institution, field) pairs and by checking whether each field's leading institution remains \#1.}

\subsection{\textcolor{red}{External validation against the Leiden Ranking}}
\label{subsec:methods-leiden}
\textcolor{red}{To assess whether $\varepsilon_2$ captures genuine research quality rather than a residual size artefact, I compare both raw h\textsubscript{2} and $\varepsilon_2$ against the CWTS Leiden Ranking Open Edition 2024~\citep{leiden2024}, an independent, OpenAlex-based institution-level bibliometric ranking based on mean normalised citation score (MNCS). I compute Spearman rank correlations $\rho(\mathrm{h}_2, \mathrm{MNCS})$ and $\rho(\varepsilon_2, \mathrm{MNCS})$ on this matched set; p-values use the standard $t$-approximation $t = \rho\sqrt{(n-2)/(1-\rho^2)}$ evaluated against a $t_{n-2}$ distribution.}

\subsection{Software and reproducibility}
All analyses were implemented in Python~3.10 (NumPy, pandas, geopandas, matplotlib, duckdb, powerlaw, scipy). The data fetching, filtering, and processing are fully reproducible from the public code repository at \url{https://github.com/rupumped/successive-h-indices}.